\section{Выводы}
Выполнив лабораторную работу №7-8 по курсу \enquote{Информационный поиск}, я
реализовал систему булева поиска над корпусом HTML-документов. Основное, что
я узнал нового:

\begin{itemize}
    \item позиционный обратный индекс на практике: для каждого терма хранятся
    не только идентификаторы документов, но и позиции вхождений, что даёт
    возможность фразового поиска без обращения к исходным документам;
    \item операции AND/OR/NOT над индексом сводятся к пересечению, объединению
    и дополнению отсортированных списков идентификаторов документов ---
    линейному по размеру списков;
    \item бинарная сериализация индекса (save/load) избавляет от
    перестроения индекса при каждом запуске: построение по корпусу из 15000
    документов занимает $\sim$35~с, загрузка из файла и поиск --- мгновенны;
    \item побайтовая обработка UTF-8 требует внимания: кириллица кодируется
    двумя байтами, и проверки диапазонов символов приходится делать по кодам
    байтов (0xD0/0xD1 и вторые байты), а не по символьным литералам;
    \item ручное управление памятью в C++ таит ловушки: в работе были найдены
    и исправлены двойное освобождение индекса (free по стековой структуре) и
    висячие указатели в демонстрационном режиме.
\end{itemize}

Для дальнейшего развития системы можно добавить:
\begin{itemize}
    \item сжатие позиционных списков (дельта-кодирование совместно с
    Variable Byte Encoding);
    \item кэширование часто используемых термов;
    \item полноценный парсер запросов с произвольной вложенностью булевых
    операторов и NEAR-поиск по окну позиций.
\end{itemize}

\pagebreak
